

Proteins carry out their diverse biological roles - from catalysis to immune recognition - by engaging in specific
physical interactions with other molecules.
These interactions are mediated by the protein's three-dimensional structure: the precise spatial arrangement of
atoms at a binding site determines which molecules the protein can interact with, and how strongly.
Predicting how a protein's structure gives rise to specific interactions - and how changes to the protein's amino-acid
sequence alter them - is therefore central to understanding disease mechanisms and designing novel therapeutics.

Physics-based approaches such as Molecular Dynamics and Rosetta provide principled, generalizable models of these
interactions, but remain severely bottlenecked by computational cost.
Machine learning has emerged as a compelling alternative, learning effective theories directly from data -
coarse-grained approximations that do not need to resolve individual inter-atomic interactions to remain predictive of
higher-level properties.
However, the central challenge in developing useful ML models for proteins lies in carefully balancing model
expressivity, training data availability, and representation choice, such that the learned shortcuts remain
physically meaningful and predictive across diverse settings.

In this thesis, I argue that this balance is best achieved by letting first-principles reasoning guide ML design choices.
The physical principle that protein function emerges from local inter-atomic interactions motivates two concrete choices:
modeling \textit{residue-centered local atomic environments} (neighborhoods), and enforcing \textit{rotational
equivariance}, ensuring that the models' predictions respect the rotational symmetry inherent to inter-atomic interactions.
These choices are combined into a unified SO(3)-equivariant neural network framework,
which I use to develop three models spanning diverse aspects of protein structure and function.

First, I introduce H-(V)AE, a rotationally equivariant (Variational) AutoEncoder that learns
compact embeddings of local atomic environments.
These embeddings encode key biophysical and structural properties and yield state-of-the-art performance on predicting
protein-ligand binding affinity.
I then present H-Packer, which uses the framework to tackle rotamer packing - the task of
predicting physically correct side-chain conformations consistent with a given backbone - achieving competitive
performance against leading physics-based and ML-based methods, and seeing broad adoption by the community.
Last, I introduce HERMES, a structure-based model pre-trained to predict amino-acid propensities
from local atomic contexts, and used in zero-shot - or after efficient fine-tuning - to predict mutational effects on
protein stability and protein-protein binding affinity, matching or exceeding state-of-the-art performance.
I also characterize HERMES' strengths, weaknesses, and biases.

Finally, I demonstrate HERMES' applicability to two distinct biological domains.
First, I show that HERMES can speed-up structure-based vaccine antigen design pileines, by scanning whole antigens
and suggesting mutations likely to stabilize viral antigens in their most immunogenic state.
Second, I leverage HERMES to address the problem of T-Cell Receptor (TCR) specificity:
predicting which peptide antigens a given TCR will recognize, a key open problem in theoretical immunology.
Without any explicit training on TCR-pMHC interaction data, HERMES successfully models the distribution of peptides
recognized by a target TCR, enabling de novo design of immunogenic peptides and achieving experimental T-cell activation
with a success rate of up to 50\%.

Taken together, the results in this thesis demonstrate that physical locality and rotational equivariance are effective
inductive biases for machine learning models of proteins, yielding a fast, principled, and broadly applicable framework
for protein modeling and design.
